\section{From computation paths to shared witnesses}\label{sec:foundations}

Suppose a proposed solution can be checked efficiently. A path through a
graph, for example, is checked by reading its vertices and checking its
successive edges. Checking one proposed path and finding a path are different
tasks. Nondeterminism gives a mathematical language for the first task
together with an existence condition: an input is accepted when at least
one legal proposal passes the check. Understanding that condition is the
starting point for both the circuit-size question and the SAT algorithm.

\subsection{A computation tree describes possibilities}

A deterministic computation has one next step at each nonhalting
configuration. A nondeterministic computation may have several legal next
steps. Starting from a fixed input, unfold these possibilities into a tree.
A root-to-leaf path is one computation, and an accepting leaf records one
successful computation. The machine accepts the input if there exists an
accepting path; it rejects if all paths reject. In the time-bounded setting
used here, every path halts within the specified bound. Its nondeterministic
running time measures the longest path, not the total number of tree nodes.

After encoding the finitely many possible transitions by a fixed number
of bits, a path can be specified by its successive choices. A deterministic
checker receives the original input and this choice record, simulates the
selected steps, rejects illegal choices, and checks that the outcome accepts.
Conversely, a machine can guess a proposed choice record and run the checker.
Thus an accepting path can be replaced by a finite \emph{certificate}, also
called a \emph{witness}. A certificate need not literally list machine
transitions: a satisfying assignment or a proposed graph path may be a
more convenient representation of the same acceptance condition.

This replacement describes existence, without promising an efficient
procedure for discovering the certificate. An acceptance proof can exhibit
one certificate. A deterministic rejection procedure must justify the
absence of every certificate, possibly by exploiting shared structure rather
than inspecting them one at a time. Nor does the computation tree provide
free physical parallelism: an implementation that assigns one processor to
every path must account for those processors, their work, and their storage.
The abstract path-length measure does not supply that hardware.

\subsection{The verifier relation and its projection}

Write $\B=\{0,1\}$. For the moment fix lengths $n$ and $m$, and let
\[
 f:\B^n\times\B^m\longrightarrow\B
\]
be a deterministic checker. The string $x$ is the \emph{ordinary input},
which specifies the question; $y$ is the proposed witness. The equation
$f(x,y)=1$ means that this particular proposal is valid. The set of valid
witnesses over $x$, called its \emph{fiber}, is
\[
 W_x=\{y\in\B^m:f(x,y)=1\}.
\]
Forgetting which witness worked leaves the Boolean function
\begin{equation}\label{eq:foundations-projection}
 g(x)=\exists y\,f(x,y)
     =\bigvee_{y\in\B^m} f(x,y)
     =\begin{cases}1,&W_x\ne\varnothing,\\0,&W_x=\varnothing.\end{cases}
\end{equation}
Here an existential assertion is identified with its Boolean truth value.
Geometrically, the accepted pairs form a relation inside
$\B^n\times\B^m$, and $g^{-1}(1)$ is its projection onto the first
coordinate. Many pairs may project to the same ordinary input. Projection
retains whether a fiber is empty and discards its size and internal shape.

We represent a checker by a \emph{Boolean circuit}: gates compute Boolean
operations and wires carry their input and output values, with no directed
feedback. A gate's \emph{fan-in} counts incoming wires; a value's
\emph{fan-out} counts its uses. The \emph{output cone} consists of the
vertices with a directed path to the designated output, including that
output. Vertices outside it can be discarded. An input is \emph{essential}
if flipping it changes the function on some assignment of the other inputs;
merely occurring in the output cone does not guarantee this property.

\begin{table}[htbp]
\centering
\begin{tabularx}{\linewidth}{@{}lX@{}}
\toprule
Object & Operational meaning \\
\midrule
$x\in\B^n$ & The instance held fixed while proposals are checked. \\
$y\in\B^m$ & One complete proposal, shared by every check that uses it. \\
$f(x,y)$ & The deterministic answer for that instance and proposal. \\
$W_x$ & All successful proposals for the fixed instance. \\
$g(x)$ & Whether there is at least one successful proposal. \\
$s$ & Number of gates in a supplied verifier circuit, when one is given. \\
$\C(g)$ & Minimum number of gates in any circuit computing $g$. \\
\bottomrule
\end{tabularx}
\caption{The same relation viewed as verification, existence, and circuit size.
The gate model is specified below; a supplied size need not be minimum.}
\label{tab:foundations-vocabulary}
\end{table}

\paragraph{A running example.}
Consider the three-gate checker
\begin{equation}\label{eq:foundations-example}
 v_1=x\lor y,\qquad v_2=y\oplus z,\qquad
 f(x,y,z)=v_1\land v_2,
\end{equation}
where $\oplus$ denotes exclusive OR: it is one exactly when its two inputs
differ. In this example $x$ is a single ordinary bit and the ordered pair
$(y,z)$ is the witness. When $x=0$, the OR gate requires $y=1$, and the
exclusive OR then requires $z=0$. When $x=1$, the OR gate already outputs
one, and either pair with $y\ne z$ works. Thus
\[
 W_0=\{(1,0)\},\qquad W_1=\{(0,1),(1,0)\}.
\]
The accepting triples, written in the order $xyz$, are precisely
$010$, $101$, and $110$. Both fibers are nonempty, so
$g(x)=\exists y,z\,f(x,y,z)=1$ for both values of $x$.
The checker uses three gates; the projected function is constant and needs
no gates under the free-constants convention. The number of witness choices
alone therefore cannot determine the difficulty of the projected function.

\subsection{Existential choices are globally consistent}

The witness is chosen once. Every occurrence of the same bit must use the
same value. A minimal failure of independent local reasoning is
\[
 \begin{gathered}
 (\exists y\in\B\;y=1)\ \land\ (\exists y\in\B\;y=0)
       \quad\text{is true},\\
 \exists y\in\B\;(y=1\land y=0)
       \quad\text{is false}.
 \end{gathered}
\]
In the left expression the quantifiers have separate scopes; renaming
their variables to $y_1$ and $y_2$ makes the independence explicit.
In the right expression one value must satisfy both conditions. In general,
existence of a common witness implies existence of separate witnesses,
but the converse fails. Independent local existence tests can admit
combinations that no actual computation realizes.

The same issue is visible inside the running circuit. Hold $x=0$ and
$z=1$. The relation $x\lor y=1$ can be satisfied with $y=1$, while
$y\oplus z=1$ can be satisfied with $y=0$. Both local tests succeed,
but no shared $y$ makes the final AND equal one. A correct algorithm may
split the two occurrences into temporary variables only if it also enforces
their equality.

Forward evaluation of a circuit never encounters a directed cycle: each
gate uses already determined predecessors. Yet the undirected pattern of
dependencies can contain a cycle. In this example the value of $y$ feeds
both $v_1$ and $v_2$, and those branches meet at the output gate. Ignoring
directions gives the cycle $y,v_1,f,v_2,y$. Such a cycle expresses shared
information, not a feedback loop in the computation. The later consistency
construction keeps this sharing explicit through equality constraints.
It is this graph of constraints, rather than the directed evaluation order
alone, that controls how many unresolved choices an elimination algorithm
must remember.

For the undirected graphs used below, a vertex's \emph{degree} counts its
incident edge ends. A graph is \emph{cubic} when every degree is three and
\emph{subcubic} when every degree is at most three. A \emph{simple} graph
has no loops or parallel edges. Intermediate constraint graphs may have
both: a loop joins a vertex to itself and contributes two to its degree,
while parallel edges join the same pair of vertices more than once.

\subsection{Existence and probability ask different questions}

Giving a witness distribution changes the question. If $Y$ is uniform on
$\B^m$, then
\[
 \Prb[f(x,Y)=1]=\frac{|W_x|}{2^m},
 \qquad
 g(x)=1\quad\Longleftrightarrow\quad\Prb[f(x,Y)=1]>0.
\]
A positive probability can be exponentially small. If a checker accepts
only the witness $0^m$, one uniform trial succeeds with probability $2^{-m}$.
For $r$ independent trials the success probability is
$1-(1-2^{-m})^r\le r2^{-m}$. Polynomially many such trials do not provide
a constant success guarantee as $m$ grows. This example concerns uniform
guessing; the accepted witness is transparent, so it is not a lower bound
against other search algorithms.

Randomized computation specifies how choices are distributed and requires
an appropriate success or error guarantee. Existential acceptance specifies
only that some legal choice succeeds. The two descriptions can use the same
tree while placing different conditions on its leaves. A good randomized
search algorithm may exploit the verifier's structure, but its guarantee
does not follow merely from the existence of one accepting path.

\subsection{From bounded certificates to polynomial verification}

To discuss a complexity class, one verifier must handle inputs of every
length. A \emph{language} is a set of finite bit strings, representing its
yes-instances. The class $\mathsf P$ consists of languages decided by a
deterministic algorithm in time polynomial in the input length. A language
belongs to $\NP$ when there are a polynomially bounded,
polynomial-time computable witness length $p(n)$ and a deterministic polynomial-time
verifier $V$ such that
\begin{equation}\label{eq:foundations-np}
 x\in L\quad\Longleftrightarrow\quad
 \exists y\in\B^{p(|x|)}\; V(x,y)=1.
\end{equation}
One may take $p$ to be an integer polynomial bound. Guessing the witness
and checking it gives a polynomial-time nondeterministic machine.
Conversely, recording the choices of a polynomial-time nondeterministic
machine gives a polynomial-length witness that can be checked by simulation.
Illegal transition codes are rejected. This proves the equivalence of the
path and verifier definitions at the level of acceptance.
The question $\mathsf P=\NP$ asks whether every language with such efficient
verification also has an efficient deterministic decision algorithm.

The fixed length in~\eqref{eq:foundations-np} is a convention that needs
care when multiplicities matter. Suppose the original verifier accepts
witnesses of any length at most $p(n)$. Encode such a witness by its length
$\ell$ in $h=\lceil\log_2(p(n)+1)\rceil$ bits, followed by $p(n)$
data bits. Require $0\le\ell\le p(n)$ and require every data bit after
position $\ell$ to be zero; run the original verifier on the first $\ell$
bits. Each original witness has exactly one valid encoding of length
$h+p(n)$. Without the suffix-zero check, a length-$\ell$ witness would
have $2^{p(n)-\ell}$ encodings. Existence would be preserved, but its
contribution to a count would change.

Similarly, a record of machine choices is not automatically in bijection
with accepting paths: unused choices after early halting or several bit
codes for one transition can multiply representations. Canonical padding
and a one-to-one transition encoding can remove these multiplicities.
Whenever counting is claimed to be preserved, the relevant bijection must
be supplied; an acceptance equivalence alone is insufficient.

For each fixed $n$, a polynomial-time verifier can be expressed by a
polynomial-size Boolean circuit $f_n(x,y)$. One elementary construction
encodes the verifier's configuration at each time step and connects a
Boolean circuit for one transition to the next. A polynomial time bound
limits both the number of steps and the tape region that can be visited,
so this finite unrolling has polynomial size. The input and witness bits
remain circuit inputs, while the machine description and time bound fix
its wiring. This construction is effective across input lengths.

Encoding an accepting computation by jointly satisfied Boolean constraints
is the classical foundation behind the SAT connection. Cook constructs a
formula satisfiable exactly when a polynomial-time nondeterministic
computation accepts, and also discusses bounded existential quantification
over deterministic polynomial-time relations \citep{Cook1971}.
Levin's foundational formulation studies universal search problems
\citep{Levin1973}. The circuit presentation here uses the same
verification principle. With $x$ fixed, asking whether $f_n(x,y)=1$ for
some $y$ is circuit satisfiability. Introducing internal gate values and
enforcing each gate relation converts this to a finite constraint problem;
once its input assignment is fixed, deterministic evaluation gives exactly
one consistent assignment to those internal gate values.

SAT abbreviates Boolean satisfiability: does some assignment make the
given formula or circuit output one? A \emph{literal} is a variable or its
negation, a \emph{clause} is an OR of literals, and a formula in
\emph{conjunctive normal form} (CNF) is an AND of clauses. In a $k$-CNF
each clause has at most $k$ literals. Gate count and clause count are
different size measures, so the SAT bounds below specify their circuit basis.

\subsection{Four computational tasks and their outputs}

For a fixed ordinary input, \emph{decision} asks whether $W_x$ is nonempty,
\emph{search} asks for one member when it is nonempty, and \emph{exact
witness counting} asks for the integer $|W_x|$. For a circuit with both
ordinary and witness inputs left variable, two totals must also be separated:
\begin{equation}\label{eq:foundations-counts}
 N_{\mathrm{pairs}}=\sum_{x\in\B^n}|W_x|
                  =\sum_{x,y}f(x,y),\qquad
 N_{\mathrm{proj}}=\sum_{x\in\B^n}[W_x\ne\varnothing]
                 =\sum_x g(x).
\end{equation}
The brackets denote an indicator, equal to one when the condition holds
and zero otherwise. The first total counts accepting input--witness pairs;
the second is \emph{projected model counting}, which counts each successful
ordinary input once. They agree when every nonempty fiber has size one,
but not in general. In the running example they are three and two.

Across input lengths, the class $\#\mathsf P$ consists of functions counting
the accepting paths of polynomial-time nondeterministic machines.
Equivalently, they count accepted fixed-length polynomially bounded
certificates for a polynomial-time verifier, using the encoding discipline
above. These are integer-valued functions, whereas $\NP$ describes decision
problems. The distinction will matter when the appendix compares existence
with exact counting.

An exact witness count decides existence by comparison with zero. It also
supports search: fix successive witness bits to a branch with a nonzero
completion count. The invariant is that the current prefix still extends
to an accepted witness, so a full prefix obtained this way is a witness.
This explains the logical reduction; preserving a particular algorithmic
exponent across these restricted instances needs the additional analysis
given later. A single decision answer, by contrast, provides no fiber size.

Counts also specify distributions. Uniformly choosing one of the three
accepting triples in the example gives $\Prb[x=1]=2/3$, whereas uniformly
choosing a successful projected input gives $\Prb[x=1]=1/2$.
In general the marginal distribution of a uniform accepting pair assigns
mass $|W_x|/N_{\mathrm{pairs}}$ to $x$. Forgetting a sampled witness
therefore favors larger fibers. The exact-counting algorithm in this
paper counts satisfying assignments to all named circuit inputs; when
these inputs are partitioned into $x$ and $y$, it computes
$N_{\mathrm{pairs}}$, without automatically computing $N_{\mathrm{proj}}$.

\subsection{Circuit existence and uniform computation}

A size bound for $g$ asserts that a circuit computing $g$ exists. A
uniform SAT algorithm receives a circuit description and computes whether
it has a satisfying assignment. These statements have different inputs
and quantifiers. At each length there could be a small circuit without
an efficient procedure for finding it. In particular, a projection that
happens to be constant has a zero-gate circuit even if recognizing that
fact from the supplied verifier is difficult.

Across lengths, $\Ppoly$ denotes languages decided by polynomial-size
circuit families, with no requirement that a single algorithm generate
the family efficiently. The verifier construction above starts with a
single polynomial-time algorithm and produces circuits effectively.
An arbitrary family of polynomial-size verifier circuits need not have
this uniformity, so its projected language should not be identified with
$\NP$ merely from the size assertion. Likewise, the polynomial size of
the verifier says little by itself about the minimum circuit size after
projection; controlling that increase is a principal question here.

For an algorithm on an explicitly supplied circuit, the input includes its
gate operations, wiring, output, and named inputs. Polynomial factors in
the running-time bounds account for this description, including inputs
that may become unused in normalization. For circuit-size bounds the
gate count is instead the resource, with the basis and free-wiring
conventions stated explicitly. Neither measure substitutes for the other.

The rest of the paper studies existential witnesses in this finite
Boolean setting. Alternation adds universal as well as existential
choices, and nondeterministic space measures memory along computations;
both lead to further questions beyond the present scope. Our route is
more specific: one accepting path gives one complete witness; fixing
that witness makes verification deterministic; shared values require
equality; and these equalities let us count consistent assignments by
combining local tables. The next construction makes that route operational.
